% CP-VI-0039
% target_chapter: VI/25 — Fibers and Geometric Fibers
% source_unit_id: IMP-DL-D27460A55B-P008
% source: Downloads/theory-of-algebraic-geometry-5.tex L5536-L6150
% solution_provenance: SOURCE_IMPLICIT_SOLUTION

\subsection*{Problem 9. Fibres of Morphisms of Affine Schemes}

\textbf{Problem.}
Describe the fibres over all points of the target space for the following
morphisms between affine schemes. In each case, the corresponding
homomorphism of rings is the obvious one. Here \(k\) denotes a field.
Which fibres are irreducible or reduced?

\textbf{General principle.}
If
\[
\varphi:A\longrightarrow B
\]
is a ring homomorphism, then the corresponding morphism of affine schemes
is
\[
\operatorname{Spec}B\longrightarrow \operatorname{Spec}A.
\]
For a point \(\mathfrak{p}\in \operatorname{Spec}A\), the fibre over
\(\mathfrak{p}\) is
\[
\operatorname{Spec}\left(B\otimes_A \kappa(\mathfrak{p})\right),
\]
where
\[
\kappa(\mathfrak{p})
=
\operatorname{Frac}(A/\mathfrak{p})
\]
is the residue field at \(\mathfrak{p}\).

\subsection{Part 1}

Consider
\[
\operatorname{Spec} k[T,U]/(TU-1)
\longrightarrow
\operatorname{Spec} k[T].
\]

Let \(\mathfrak{p}\in \operatorname{Spec}k[T]\), and write
\[
a=\overline{T}\in \kappa(\mathfrak{p}).
\]
The fibre is
\[
\operatorname{Spec}
\left(
\kappa(\mathfrak{p})[U]/(aU-1)
\right).
\]

If \(\mathfrak{p}=(T)\), then \(a=0\), so the fibre is
\[
\operatorname{Spec}\left(k[U]/(-1)\right)
=
\operatorname{Spec}(0)
=
\varnothing.
\]

If \(\mathfrak{p}\neq (T)\), then \(a\neq 0\) in \(\kappa(\mathfrak{p})\),
so
\[
\kappa(\mathfrak{p})[U]/(aU-1)
\cong
\kappa(\mathfrak{p}).
\]
Hence the fibre is a single reduced point:
\[
\operatorname{Spec}\kappa(\mathfrak{p}).
\]

Thus:
\[
\begin{array}{c|c|c|c}
	\text{Point of } \operatorname{Spec}k[T]
	&
	\text{Fibre}
	&
	\text{Reduced?}
	&
	\text{Irreducible?}
	\\
	\hline
	(T)
	&
	\varnothing
	&
	\text{yes}
	&
	\text{no, by the usual convention}
	\\
	\mathfrak{p}\neq (T)
	&
	\operatorname{Spec}\kappa(\mathfrak{p})
	&
	\text{yes}
	&
	\text{yes}
\end{array}
\]

Geometrically, this is the morphism
\[
\mathbb{A}^1_k\setminus \{0\}\longrightarrow \mathbb{A}^1_k,
\]
so the fibre over \(T=0\) is empty and every other fibre is one point.

\subsection{Part 2}

Consider
\[
\operatorname{Spec} k[T,U]/(T^2-U^2)
\longrightarrow
\operatorname{Spec} k[T].
\]

Again let
\[
a=\overline{T}\in \kappa(\mathfrak{p}).
\]
The fibre is
\[
\operatorname{Spec}
\left(
\kappa(\mathfrak{p})[U]/(a^2-U^2)
\right).
\]
Equivalently,
\[
a^2-U^2=(a-U)(a+U).
\]

If \(\mathfrak{p}=(T)\), then \(a=0\), so the fibre is
\[
\operatorname{Spec}\left(k[U]/(U^2)\right).
\]
This is one nonreduced point. Hence it is irreducible but not reduced.

Now suppose \(\mathfrak{p}\neq (T)\), so \(a\neq 0\).

If \(\operatorname{char}k\neq 2\), then \(a\neq -a\), and
\[
a^2-U^2=(a-U)(a+U)
\]
has two distinct linear factors. Therefore
\[
\kappa(\mathfrak{p})[U]/(a^2-U^2)
\cong
\kappa(\mathfrak{p})\times \kappa(\mathfrak{p}).
\]
The fibre consists of two reduced points. It is reduced but not
irreducible.

If \(\operatorname{char}k=2\), then \(a=-a\), and
\[
a^2-U^2=(U-a)^2.
\]
Hence the fibre is
\[
\operatorname{Spec}
\left(
\kappa(\mathfrak{p})[U]/((U-a)^2)
\right).
\]
This is one nonreduced point. It is irreducible but not reduced.

Thus, if \(\operatorname{char}k\neq 2\),
\[
\begin{array}{c|c|c|c}
	\text{Point}
	&
	\text{Fibre}
	&
	\text{Reduced?}
	&
	\text{Irreducible?}
	\\
	\hline
	(T)
	&
	\operatorname{Spec}k[U]/(U^2)
	&
	\text{no}
	&
	\text{yes}
	\\
	\mathfrak{p}\neq (T)
	&
	\text{two reduced points}
	&
	\text{yes}
	&
	\text{no}
\end{array}
\]

If \(\operatorname{char}k=2\), every fibre is a doubled point, hence
irreducible but not reduced.

\subsection{Part 3}

Consider
\[
\operatorname{Spec}
k[T,U,V,W]/
\bigl((U+T)W,\ (U+T)(U^3+U^2+UV^2-V^2)\bigr)
\longrightarrow
\operatorname{Spec}k[T].
\]

Let
\[
a=\overline{T}\in \kappa(\mathfrak{p}).
\]
The fibre is
\[
\operatorname{Spec}
\frac{\kappa(\mathfrak{p})[U,V,W]}
{
	((U+a)W,\ (U+a)(U^3+U^2+UV^2-V^2))
}.
\]

Set
\[
F(U,V)=U^3+U^2+UV^2-V^2.
\]
Then the fibre is cut out by
\[
(U+a)W=0,
\qquad
(U+a)F(U,V)=0.
\]

Set-theoretically, this fibre is the union
\[
\{U=-a\}
\cup
\{W=0,\ F(U,V)=0\}.
\]
The first part is an affine plane:
\[
\operatorname{Spec}\kappa(\mathfrak{p})[V,W].
\]
The second part lies inside the plane \(W=0\) and is the curve
\[
F(U,V)=0.
\]

Therefore every fibre is reducible: it contains the plane \(U=-a\) and
also the residual part \(W=0,\ F(U,V)=0\).

Now discuss reducedness.

If \(\operatorname{char}k\neq 2\), then \(F(U,V)\) is square-free, and
the factor \(U+a\) is not a repeated factor of the defining equations.
In this case the ideal
\[
((U+a)W,\ (U+a)F)
\]
is radical. Hence every fibre is reduced.

If \(\operatorname{char}k=2\), then
\[
F(U,V)
=
U^3+U^2+UV^2+V^2
=
(U+1)(U+V)^2.
\]
Thus the equations contain a repeated factor. In this case every fibre is
nonreduced.

Hence:
\[
\begin{array}{c|c|c}
	\text{Characteristic of }k
	&
	\text{Reducedness of fibres}
	&
	\text{Irreducibility of fibres}
	\\
	\hline
	\operatorname{char}k\neq 2
	&
	\text{all fibres reduced}
	&
	\text{all fibres reducible}
	\\
	\operatorname{char}k=2
	&
	\text{all fibres nonreduced}
	&
	\text{all fibres reducible}
\end{array}
\]

\subsection{Part 4}

Consider
\[
\operatorname{Spec}\mathbb{Z}[T]
\longrightarrow
\operatorname{Spec}\mathbb{Z}.
\]

The points of \(\operatorname{Spec}\mathbb{Z}\) are the generic point
\[
(0)
\]
and the closed points
\[
(p),
\]
where \(p\) is a prime number.

Over the generic point \((0)\), the residue field is
\[
\kappa((0))=\mathbb{Q}.
\]
The fibre is
\[
\operatorname{Spec}\left(\mathbb{Z}[T]\otimes_{\mathbb{Z}}\mathbb{Q}\right)
=
\operatorname{Spec}\mathbb{Q}[T]
=
\mathbb{A}^1_{\mathbb{Q}}.
\]

Over a closed point \((p)\), the residue field is
\[
\kappa((p))=\mathbb{F}_p.
\]
The fibre is
\[
\operatorname{Spec}\left(\mathbb{Z}[T]\otimes_{\mathbb{Z}}\mathbb{F}_p\right)
=
\operatorname{Spec}\mathbb{F}_p[T]
=
\mathbb{A}^1_{\mathbb{F}_p}.
\]

Thus:
\[
\begin{array}{c|c|c|c}
	\text{Point of }\operatorname{Spec}\mathbb{Z}
	&
	\text{Fibre}
	&
	\text{Reduced?}
	&
	\text{Irreducible?}
	\\
	\hline
	(0)
	&
	\operatorname{Spec}\mathbb{Q}[T]
	&
	\text{yes}
	&
	\text{yes}
	\\
	(p)
	&
	\operatorname{Spec}\mathbb{F}_p[T]
	&
	\text{yes}
	&
	\text{yes}
\end{array}
\]

So every fibre is an affine line over a field.

\subsection{Part 5}

Consider
\[
\operatorname{Spec}\mathbb{Z}[T]/(T^2+1)
\longrightarrow
\operatorname{Spec}\mathbb{Z}.
\]

Over the generic point \((0)\), the fibre is
\[
\operatorname{Spec}
\left(
\mathbb{Q}[T]/(T^2+1)
\right).
\]
Since \(T^2+1\) is irreducible over \(\mathbb{Q}\), we get
\[
\mathbb{Q}[T]/(T^2+1)\cong \mathbb{Q}(i).
\]
Thus the generic fibre is
\[
\operatorname{Spec}\mathbb{Q}(i),
\]
a single reduced point.

Now let \(p\) be a prime number. The fibre over \((p)\) is
\[
\operatorname{Spec}
\left(
\mathbb{F}_p[T]/(T^2+1)
\right).
\]

For \(p=2\), we have
\[
T^2+1=(T+1)^2
\]
in \(\mathbb{F}_2[T]\). Hence the fibre is
\[
\operatorname{Spec}\left(\mathbb{F}_2[T]/((T+1)^2)\right),
\]
a nonreduced one-point scheme. It is irreducible but not reduced.

For \(p\) odd, the polynomial \(T^2+1\) splits over \(\mathbb{F}_p\) if
and only if \(-1\) is a square modulo \(p\). By the classical number
theoretic criterion,
\[
-1 \text{ is a square modulo }p
\iff
p\equiv 1 \pmod 4.
\]

Therefore:

If
\[
p\equiv 1 \pmod 4,
\]
then
\[
T^2+1=(T-a)(T+a)
\]
for some \(a\in \mathbb{F}_p\), and the fibre is
\[
\operatorname{Spec}(\mathbb{F}_p\times \mathbb{F}_p).
\]
It consists of two reduced points. Hence it is reduced but not
irreducible.

If
\[
p\equiv 3 \pmod 4,
\]
then \(T^2+1\) is irreducible over \(\mathbb{F}_p\), so
\[
\mathbb{F}_p[T]/(T^2+1)\cong \mathbb{F}_{p^2}.
\]
The fibre is one reduced point:
\[
\operatorname{Spec}\mathbb{F}_{p^2}.
\]

Thus:
\[
\begin{array}{c|c|c|c}
	\text{Point of }\operatorname{Spec}\mathbb{Z}
	&
	\text{Fibre}
	&
	\text{Reduced?}
	&
	\text{Irreducible?}
	\\
	\hline
	(0)
	&
	\operatorname{Spec}\mathbb{Q}(i)
	&
	\text{yes}
	&
	\text{yes}
	\\
	(2)
	&
	\operatorname{Spec}\mathbb{F}_2[T]/((T+1)^2)
	&
	\text{no}
	&
	\text{yes}
	\\
	(p),\ p\equiv 1\pmod 4
	&
	\text{two reduced points}
	&
	\text{yes}
	&
	\text{no}
	\\
	(p),\ p\equiv 3\pmod 4
	&
	\operatorname{Spec}\mathbb{F}_{p^2}
	&
	\text{yes}
	&
	\text{yes}
\end{array}
\]

From the point of view of algebraic number theory, this is the behaviour
of rational primes in the Gaussian integers \(\mathbb{Z}[i]\):
\[
\begin{array}{c|c}
	\text{Prime }p
	&
	\text{Behaviour in }\mathbb{Z}[i]
	\\
	\hline
	p=2
	&
	\text{ramified}
	\\
	p\equiv 1\pmod 4
	&
	\text{split}
	\\
	p\equiv 3\pmod 4
	&
	\text{inert}
\end{array}
\]

\subsection{Part 6}

Consider
\[
\operatorname{Spec}\mathbb{C}
\longrightarrow
\operatorname{Spec}\mathbb{Z}.
\]

The corresponding ring homomorphism is
\[
\mathbb{Z}\longrightarrow \mathbb{C}.
\]

The scheme \(\operatorname{Spec}\mathbb{C}\) has only one point, namely
the zero ideal \((0)\). Its inverse image in \(\operatorname{Spec}\mathbb{Z}\)
is
\[
(0)\subseteq \mathbb{Z}.
\]
Thus the image of the morphism is only the generic point of
\(\operatorname{Spec}\mathbb{Z}\).

The fibre over \((0)\) is
\[
\operatorname{Spec}
\left(
\mathbb{C}\otimes_{\mathbb{Z}}\mathbb{Q}
\right)
=
\operatorname{Spec}\mathbb{C}.
\]
This is one reduced point.

For a closed point \((p)\), the fibre is
\[
\operatorname{Spec}
\left(
\mathbb{C}\otimes_{\mathbb{Z}}\mathbb{F}_p
\right).
\]
But
\[
\mathbb{C}\otimes_{\mathbb{Z}}\mathbb{F}_p
\cong
\mathbb{C}/p\mathbb{C}.
\]
Since \(p\) is invertible in \(\mathbb{C}\), we have
\[
p\mathbb{C}=\mathbb{C}.
\]
Therefore
\[
\mathbb{C}/p\mathbb{C}=0.
\]
Hence the fibre over \((p)\) is empty.

Thus:
\[
\begin{array}{c|c|c|c}
	\text{Point of }\operatorname{Spec}\mathbb{Z}
	&
	\text{Fibre}
	&
	\text{Reduced?}
	&
	\text{Irreducible?}
	\\
	\hline
	(0)
	&
	\operatorname{Spec}\mathbb{C}
	&
	\text{yes}
	&
	\text{yes}
	\\
	(p)
	&
	\varnothing
	&
	\text{yes}
	&
	\text{no, by the usual convention}
\end{array}
\]
